\section{Background}
\label{sec:background}

This section is deliberately short. Its job is to define the terms the rest of the paper uses and
to sketch the system's shape, not to re-argue the literature placement already made in
Section~\ref{sec:related-work}.

\subsection{Recirculating aquaculture systems}

A recirculating aquaculture system (RAS) treats and reuses water within a closed loop rather than
drawing continuously from an open source, which makes water chemistry both more controllable and
more fragile: a single mechanical or biological failure propagates through the whole tank volume
rather than being diluted by fresh inflow. Dissolved oxygen (DO) is the parameter this paper treats
as the sharpest deadline. Reported lethal thresholds for farmed fish vary by species and water
temperature, roughly in the 4.1--7.0~mg/L range \citep{lindholm2023waterquality,
agmrc2026waterquality}; for the sturgeon stock this system monitors, the operative floor is
6~mg/L, and a sustained drop below it can kill stock within hours, not days. That timescale is the
reason a monitoring system for this domain cannot tolerate a silent failure mode: there is no
grace period in which a human operator is expected to notice and intervene before an LLM's
unnoticed mistake becomes a physical loss.

\subsection{The S.U.R.E. system, at a glance}

S.U.R.E. is an implemented welfare-monitoring stack, designed to run continuously, built around four
communicating services (Figure references and full architectural detail follow in
Section~\ref{sec:system-architecture}; this is only an orientation map). A vision service runs
YOLOv11s object detection with ByteTrack \citep{zhang2022bytetrack} multi-object tracking over a
live or recorded camera feed at roughly 15 frames per second, producing a per-frame fish count and
track history. A sensor stream (dissolved oxygen, temperature, pH, and related water-chemistry
readings) arrives on a two-second cycle. A backend service fuses both streams, maintains a rolling
history, and computes a decision. That decision has two independent sources: a locally hosted LLM
(AQUA-1B, a fine-tuned Gemma-3-1B variant with a LoRA domain adapter) that reasons over the fused
sensor-and-vision state and produces a natural-language explanation alongside a status label, and a
deterministic rule engine that computes the same status label directly from the sensor and vision
thresholds, with no dependency on the model at all. A frontend dashboard renders the fused result.
A retrieval-augmented generation (RAG) layer supplies the LLM with grounded passages from eight
domain-specific knowledge documents when it needs supporting context, and an MLOps layer tracks
model registry history and watches the vision model's own output distribution for drift.

\subsection{Definitions used throughout}

\paragraph{Severity levels.} The system recognizes three ordered severity levels --
\texttt{ok} $<$ \texttt{warning} $<$ \texttt{critical} -- forming a finite total order. Every
decision, whether produced by the LLM or the rule engine, resolves to exactly one of these three
labels; there is no partial-severity or confidence-weighted intermediate state exposed to either
side of the decision fusion described next.

\paragraph{The rule engine.} ``The rule engine'' refers throughout this paper to
\texttt{backend/rules.py}, a single Python module with no dependency on the model-serving stack,
the web framework, or any request/response object -- it is pure functions over sensor and vision
values. Both the system's decision path and the offline evaluation harness import this exact
module, a property Section~\ref{sec:results} verifies directly rather than assumes. The rule
engine's output is treated as ground truth for what the system's severity \emph{should} be at any
given moment; the LLM's job is to reason about and explain that state, never to override it
downward.

\paragraph{Decision fusion and escalation.} When the LLM's proposed status disagrees with the rule
engine's computed status, the system never lets the LLM's status stand if it is less severe than
the rule engine's. This escalation-only behavior is stated formally as
Formalization~\ref{form:monotonicity} in Section~\ref{sec:system-architecture}.

\paragraph{Retrieval-augmented generation (RAG).} RAG here means the specific pipeline that embeds
incoming queries and eight source documents with the \texttt{e5-small} model, retrieves
similarity-ranked passage chunks from a \texttt{pgvector} store, and supplies retrieved passages to
the LLM as grounding context, subject to a minimum-similarity threshold below which nothing is
retrieved at all. It is explicitly an enhancement to the LLM's reasoning, not a dependency the
system requires to keep functioning -- if the vector store is unreachable, the system continues to
operate on the rule engine alone.

\paragraph{Population Stability Index (PSI).} PSI is the divergence measure used here to detect
drift in the vision model's own detection-confidence output distribution, computed by
comparing a binned reference distribution (captured when the reference window was first
established) against a binned current
window using $\mathrm{PSI} = \sum_i (c_i - r_i)\ln(c_i / r_i)$ over bins $i$ with reference and
current mass $r_i, c_i$. The binning scheme used to build those bins is itself the object of one of
this paper's own MLOps findings (Section~\ref{sec:results}); PSI's sensitivity is not independent
of how its bins are drawn, and that dependency, not merely PSI's value, is what the paper's MLOps
result is actually about.
